\chapter{二十五讲闭卷提取表}

先遮住右两栏。任选一讲，用一分钟列条件和例外，再用一句反例说明其中一项为何必要。
两轮后把相邻的规则放在一起做对比，依次核对左栏的讲次。右栏是变化的输入，须自己说明变化后的结论。

\begingroup
\renewcommand{\arraystretch}{1.3}
\begin{longtable}{@{}p{.12\linewidth}p{.42\linewidth}p{.40\linewidth}@{}}
\toprule
讲次 & 闭卷提取顺序 & 最小变式\\
\midrule
\endfirsthead
\toprule
讲次 & 闭卷提取顺序 & 最小变式\\
\midrule
\endhead
\bottomrule
\endfoot
\hyperref[lecture:01]{第一讲} & 解释边界、三原则、适用法域与时间版本。 & 从语词可能范围内改为范围外，或把适用法定刑幅度调换。\\
\hyperref[lecture:02]{第二讲} & 写大前提和小前提，按要件列事实，再查违法与责任。 & 缺一个条件，或给出已经成立的阻却事由。\\
\hyperref[lecture:03]{第三讲} & 分清行为对象和所侵害法益，核对规范保护方向。 & 行为对象相同，改为不同来源的侵害方式。\\
\hyperref[lecture:04]{第四讲} & 核对自然人身份、单位犯罪法定性、单位归属和处罚主体。 & 只把行为改为成员盗用单位名义、违法所得归个人。\\
\hyperref[lecture:05]{第五讲} & 作为先找实行危险。不作为依次写义务、能做、未做及等价或法定条件。 & 有救助义务，事发时却无现实救助能力。\\
\hyperref[lecture:06]{第六讲} & 危险与实害分列，结果加重先查基本行为，再查结果、归责、罪过。 & 重结果相同，改成与基本行为无关联的来源。\\
\hyperref[lecture:07]{第七讲} & 固定具体结果，找事实条件，核对风险实现，再处理介入和不明原因。 & 保持死亡，改时间、机制或独立介入因素。\\
\hyperref[lecture:08]{第八讲} & 先写知道什么，再写希望、放任、轻信或疏忽。错误题分对象、打击与因果经过。 & 从放任改成有具体根据的轻信，或只变实际对象。\\
\hyperref[lecture:09]{第九讲} & 防卫查基础与双门槛，避险查不得已和利益限度，承诺查权限、时间和范围。 & 重伤相同，改动是否明显超限。\\
\hyperref[lecture:10]{第十讲} & 先行为时年龄，再责任能力、违法性认识可能性、期待可能性。 & 十三周岁，给定死亡与恶劣情节，只改有无最高检核准。\\
\hyperref[lecture:11]{第十一讲} & 在时间线上标着手、停止原因与完成节点，再看中止有效性。 & 同样没得逞，改“能继续而主动停止”为“被迫停止”。\\
\hyperref[lecture:12]{第十二讲} & 定共同范围、角色、参加时间，逐人查故意、责任、过限和处罚。 & 只把帮助时间从实行中改到既遂结束后且无事前通谋。\\
\hyperref[lecture:13]{第十三讲} & 分行为、列候选罪、查特别规则、定处罚个数。 & 候选罪相同，只改变是否由两个独立行为产生。\\
\hyperref[lecture:14]{第十四讲} & 先标刑种，再记刑期、起算或适用范围，区分附加刑和非刑罚措施。 & 把罚金改成追缴违法所得，问法律性质是否相同。\\
\hyperref[lecture:15]{第十五讲} & 确定法定刑档，逐项找量刑情节，再做并罚、缓刑。 & 把原判前漏罪改成原判后新罪，交换运算次序。\\
\hyperref[lecture:16]{第十六讲} & 减刑查实际执行下限，假释查准入与禁止，时效查法定最高刑和起算例外。 & 最高刑从九年改为十年，或增加立案后逃避侦查。\\
\hyperref[lecture:17]{第十七讲} & 定罪名、拆罪状、找所引规范，删除条款检验注意规定与拟制。 & 删除指定处理条款，问普通构成能否独立推出同一结论。\\
\hyperref[lecture:18]{第十八讲} & 按生命身体、性权利、人身自由、名誉人身权、婚姻家庭选路径。 & 后果相同，只改故意指向，或只改非法控制人的目的。\\
\hyperref[lecture:19]{第十九讲} & 查原占有、转移机制、取得节点、数额口径、法定转化和罪数。 & 保留欺骗，只改为被害人未交付且行为人自行取走。\\
\hyperref[lecture:20]{第二十讲} & 查公共风险的范围，再区分危险状态、事故结果与各特殊对象。 & 实际一人受害，改为风险可扩散至不特定人。\\
\hyperref[lecture:21]{第二十一讲} & 先找前置规范，再区分主体、对象、行为、金额或危险指标。 & 把销售金额换成未销售货值，问可否用同一门槛。\\
\hyperref[lecture:22]{第二十二讲} & 先找公共秩序等具体制度，再找特别身份、行为和情节。 & 上游人员未被判刑，但赃物来源另有证据。\\
\hyperref[lecture:23]{第二十三讲} & 查身份、公款去向、占有目的、本人或他人职务路径、利益及情节。 & 把直接职务制约改成借地位请无制约关系的他人。\\
\hyperref[lecture:24]{第二十四讲} & 定具体职权或职责，区分滥用与失职，查损失、归责及特别罪名。 & 同样未履职，改为故意不履行职务。\\
\hyperref[lecture:25]{第二十五讲} & 对照间谍活动与为境外获取、提供秘密情报的对象和行为分支。 & 把合格秘密对象改成普通公开资料，问哪条路径被切断。\\
\end{longtable}
\endgroup
